\documentclass[11pt]{article}
\usepackage[a4paper,margin=1in]{geometry}
\usepackage[T1]{fontenc}
\usepackage{lmodern,microtype}
\usepackage{amsmath,amssymb,amsthm,booktabs}
\usepackage[colorlinks=true,allcolors=blue]{hyperref}
\newtheorem{theorem}{Theorem}[section]
\newtheorem{proposition}[theorem]{Proposition}
\newtheorem{corollary}[theorem]{Corollary}
\theoremstyle{remark}\newtheorem{remark}[theorem]{Remark}
\title{All-Period Integral Monodromy and Algebraic-Unit\\Multipliers for the Area-Preserving H\'enon Map}
\author{Liang Wang$^{*1}$\\
\small $^{1}$School of Artificial Intelligence and Automation, Huazhong University of Science and Technology\\
\small Wuhan 430074, P.R. China\\
\small $^*$Corresponding author}
\date{Preprint, August 2026}
\hypersetup{pdfauthor={Liang Wang},pdftitle={All-Period Integral Monodromy for the Henon Map}}
\begin{document}\maketitle
\begin{abstract}
Pressure normalization gives a prime-orbit counting law for a certified
H\'enon suspension, leaving the arithmetic type of its orbit labels as the
next gate.  We prove an all-period algebraic theorem for the underlying
autonomous map $H_6(q,p)=(1-6q^2-p,q)$.  After the canonical scaling
$x_i=6q_i$, the period-$n$ fixed algebra is defined by monic cyclic
quadratics and is finite free over $\mathbb Z$ of rank $2^n$.  Every
chronological derivative step lies in $SL_2$ over this integral algebra.
Therefore every periodic monodromy trace is an algebraic integer and every
multiplier is an algebraic unit.  A separate exact implementation verifies
the Gr\"obner and trace ledgers through period 10 and recovers the inherited
degree-four fixed-point multiplier polynomial.  The theorem supplies genuine
all-period arithmetic structure, but it also rules out treating a raw
multiplier or its field norm as a rational prime.  It does not classify the
nonalgebraic pressure power $|\Lambda|^{h_*}$.
\end{abstract}

\section{Integral cyclic coordinates}
The H\'enon recurrence is
\[
q_{i+1}=1-6q_i^2-q_{i-1}.
\]
Set $x_i=6q_i$.  Then
\begin{equation}\label{eq:scaled}
x_{i+1}=6-x_i^2-x_{i-1}.
\end{equation}
For period $n$, indices are read cyclically.  The neighbor multiset must be
retained at low period, giving
\[
n=1:\quad x_0^2+2x_0-6=0,
\]
and
\[
n=2:\quad x_0^2+2x_1-6=x_1^2+2x_0-6=0.
\]

\section{Finite-free fixed algebra}
Let
\[
\mathcal A_n=\mathbb Z[x_0,\ldots,x_{n-1}]/(f_0,\ldots,f_{n-1}),
\quad
f_i=x_i^2+x_{i-1}+x_{i+1}-6,
\]
with the exceptional neighbor multiplicities above.

\begin{theorem}[Integral all-period fixed algebra]\label{thm:free}
For every $n\ge1$, $\mathcal A_n$ is finite free over $\mathbb Z$ of rank
$2^n$, with basis
\[
x_0^{\epsilon_0}\cdots x_{n-1}^{\epsilon_{n-1}},
\qquad \epsilon_i\in\{0,1\}.
\]
In particular every coordinate value at every geometric periodic point is an
algebraic integer.
\end{theorem}
\begin{proof}
Under any degree-compatible monomial order, $f_i$ is monic with leading
monomial $x_i^2$.  The leading monomials are pairwise coprime, so Buchberger's
product criterion makes the displayed equations a monic Gr\"obner basis over
$\mathbb Z$.  The standard monomials are exactly the square-free monomials,
which proves finite freeness and the rank.  Specializing a finite integral
$\mathbb Z$-algebra at a geometric point sends each $x_i$ to an algebraic
integer.
\end{proof}

This strengthens the earlier finite-flat statement over
$\mathbb Z[A,A^{-1}]$ at the integer specialization: the scaled $A=6$ fibre
needs no denominator localization.

\section{Chronological monodromy}
In either the original $(q_i,q_{i-1})$ coordinates or the linearly scaled
coordinates, the derivative step is conjugate to
\begin{equation}\label{eq:J}
J_i=\begin{pmatrix}-2x_i&-1\\1&0\end{pmatrix},
\qquad \det J_i=1.
\end{equation}
For a period-$n$ point, let
\[
M_\gamma=J_{n-1}\cdots J_0,
\qquad t_\gamma=\operatorname{tr}M_\gamma.
\]
Later times act on the left.

\begin{theorem}[Algebraic-unit multiplier theorem]\label{thm:unit}
For every geometric periodic point of $H_6$ at every period,
$t_\gamma$ is an algebraic integer and both eigenvalues of $M_\gamma$ are
algebraic units.  Their characteristic polynomial is
\[
X^2-t_\gamma X+1.
\]
\end{theorem}
\begin{proof}
Equation~\eqref{eq:J} and Theorem~\ref{thm:free} show that every entry of
$M_\gamma$, hence its trace, is an algebraic integer.  An eigenvalue
$\lambda$ is integral over the ring of algebraic integers by the displayed
monic polynomial, hence integral over $\mathbb Z$.  Since
$\det M_\gamma=1$, the other eigenvalue is $\lambda^{-1}$ and is also an
algebraic integer.  Thus $\lambda$ is a unit.
\end{proof}

The theorem includes nonreal and multiple geometric points scheme by scheme;
reducedness is not needed for the integrality statement.  On the certified
real hyperbolic survivor, the unstable multiplier is the root of modulus
greater than one.

\section{Exact sentinels}
The code builds the Gr\"obner basis and chronological trace independently for
$1\le n\le10$.  The first reduced traces are
\[
-2x_0,
\qquad
4x_0x_1-2,
\qquad
-8x_0x_1x_2+2(x_0+x_1+x_2).
\]
All coefficients are integers and all monomials lie in the square-free
basis.  At period one, $x_0=-1\pm\sqrt7$ and the two trace-conjugate
characteristic polynomials multiply to
\[
X^4-4X^3-22X^2-4X+1,
\]
recovering the exact non-lattice witness from the instability-roof project.

\section{Arithmetic consequences and boundary}
If an algebraic unit lies in $\mathbb Q$, it is $\pm1$.  Thus no hyperbolic
raw multiplier can itself be a rational prime.  Its field norm is also
$\pm1$.  Trace and $\det(I-M)$ are algebraic integers but are not
multiplicative under repetition, so they cannot simply replace the Euler
label without a new compiler.

The pressure-normalized label $|\Lambda|^{h_*}$ is deliberately outside this
corollary.  The pressure root is a real dynamical invariant whose algebraic
type is unknown; no valid transcendence conclusion is available here.

\section{Evaluator verdict}
The all-period unit theorem is positive arithmetic structure at A1/A3, but
it constructs no prime correspondence or determinant.  The strict tuple is
\[
(A1_{\rm WEAK},A2_{\rm FAIL},A3_{\rm PARTIAL},A4_{\rm FORMAL}),
\]
with overall \texttt{ROUTE\_A\_EXPLORATORY}.  Route B is not testable.

\section{Conclusion}
The H\'enon instability multipliers are not arbitrary reals: after the
correct integral scaling, they form an all-period family of algebraic units.
This both strengthens the arithmetic content of the pressure-normalized road
and closes its most naive raw-prime interpretation.  The next step is to
classify which scalar functions of a unit can preserve exact repetition and
still produce rational prime labels.

\begin{thebibliography}{9}
\bibitem{Buchberger}
B. Buchberger, An algorithmic criterion for the solvability of algebraic
systems of equations, \emph{Aequationes Math.} 4 (1970), 374--383.
\bibitem{Silverman}
J. H. Silverman, The arithmetic of dynamical systems, Springer, 2007.
\end{thebibliography}
\end{document}
